\documentclass{exam}

\newif\ifA
\newif\ifkey
\newif\ifprintselected   % required by exam_config's \ansbox

%%%%%%%%%%%%%%%%%%%%%%%%%%%%%%%%%%%%%%%%%%%%%%%%%%%%%%
%%% SET THESE IF VARIABLES
%%%%%%%%%%%%%%%%%%%%%%%%%%%%%%%%%%%%%%%%%%%%%%%%%%%%%%
% Lab name -- labs are named by the lesson they belong to, never by a
% sequential lab number (Neil, 2026-10-12): this is the L5 Lab.
\def\labLesson{5}

% is this the key or not?
\ifdefined\iskey
	\ifnum\iskey=1
		\keytrue
	\else
		\keyfalse
	\fi
\else
	%% Pick one manually
	\keyfalse % if this is NOT the key
%	\keytrue     % if this IS the key
\fi

% set up the header and footer - change for every term
\ifdefined\thisTermIn
	\def\thisterm{\thisTermIn}
\else
	\def\thisterm{Spring 2027}
\fi
%
\def\thisclass{ECE 448}

%%%%%%%%%%%%%%%%%%%%%%%%%%%%%%%%%%%%%%%%%%%%%%%%%%%%%%
%%% END SET THESE IF VARIABLES
%%%%%%%%%%%%%%%%%%%%%%%%%%%%%%%%%%%%%%%%%%%%%%%%%%%%%%
% packages and shortcuts
\input{myPackages_gr}
\usepackage[hidelinks, linkcolor=red]{hyperref}
\input{myShortcuts}
\setlength{\parindent}{0pt}
\setlength{\parskip}{6pt}

%% exam class customization
\input{exam_config}
% ECE 448 document fonts.
%
% Body text is Barlow (SIL OFL; the faces and the license are in
% latex/fonts/Barlow/). Barlow is not in TeX Live, so it is loaded as an
% OpenType face through fontspec -- which means the document must be built
% with lualatex, not pdflatex. build_practice.sh and build_lab.sh both do.
%
% Compiling with pdflatex anyway is not an error: the guard below falls back
% to the LaTeX default so the source still compiles under any engine and can
% be syntax-checked. The resulting PDF just is not the one we ship.
%
% Paths are relative to latex/, which is where both build scripts run
% (each does cd "$(dirname "$0")"). Loading by path rather than by installed
% family name means no system font install and no fontconfig, so a fresh
% clone builds a byte-comparable PDF anywhere.
%
% MATH IS DELIBERATELY LEFT IN COMPUTER MODERN. Barlow has no math companion,
% and CM math under sans body text is already the house pairing: the decks set
% words in Source Sans Pro and let MathJax render symbols in its CM-derived
% face. Moving math to a sans as well would mean unicode-math plus a sans math
% font (Fira Math being the plausible one) and re-checking every equation in
% the course for symbol coverage. That is a separate decision, not a side
% effect of this one.

\usepackage{iftex}
\ifLuaTeX
	\usepackage{fontspec}
	\defaultfontfeatures{Ligatures=TeX}
	\setmainfont{Barlow}[
		Path           = fonts/Barlow/,
		Extension      = .ttf,
		UprightFont    = *-Regular,
		ItalicFont     = *-Italic,
		BoldFont       = *-Bold,
		BoldItalicFont = *-BoldItalic,
	]
	\setsansfont{Barlow}[
		Path           = fonts/Barlow/,
		Extension      = .ttf,
		UprightFont    = *-Regular,
		ItalicFont     = *-Italic,
		BoldFont       = *-SemiBold,
		BoldItalicFont = *-BoldItalic,
	]
	% microtype earns its place here specifically because of Barlow. Its
	% glyphs are narrower than Computer Modern's, so justified prose sets
	% tighter and TeX has to squeeze interword space to fit; two paragraphs
	% in the L02/L03 sets went overfull on the first Barlow build. Protrusion
	% and font expansion absorb that -- both are lualatex features, so this
	% belongs inside the guard with everything else.
	\usepackage{microtype}

	% Mono stays Latin Modern: Barlow ships no monospaced face, and verbatim
	% and code need one that actually is fixed-pitch.
\fi
   % Barlow body text (needs lualatex)
\printselectedfalse      % full worked solutions in the key, not answers-only
\CorrectChoiceEmphasis{\color{red}}
\SolutionEmphasis{\color{red}}
\renewcommand\questionlabel{\textbf{\thequestion.}}
\renewcommand{\thepartno}{\alph{partno}}
\renewcommand{\partlabel}{(\thepartno)}

% Fillable table cell: blank in the student copy, the expected value in red on
% the KEY. Fixed-width centered columns keep the blank cells writable.
\newcolumntype{K}[1]{>{\centering\arraybackslash}p{#1}}
\newcommand{\kv}[1]{\ifkey{\color{red}#1}\fi\rule[-0.10in]{0pt}{0.34in}}

% print the answers if this is the key
\ifkey
	\printanswers % if this IS the key
\else
	\noprintanswers  %if this is NOT the key
\fi

% print the header and footer
% need to print key on top if this is the key
\ifkey
	\firstpageheader{}{\color{red}{ KEY}}{}
	\runningheader{\thisclass\hspace{1pt} -- \thisterm }{L\labLesson\ Lab\hspace{1pt}\color{red}{ KEY}}{\thepage\hspace{1pt} of\hspace{1pt} \numpages}
	\firstpagefooter{}{}{}
	\runningfooter{}{}{}
\else
	\firstpageheader{}{}{}
	\runningheader{\thisclass \hspace{1pt} -- \thisterm }{L\labLesson\ Lab\hspace{1pt}}{\thepage\hspace{1pt} of\hspace{1pt} \numpages}
	\firstpagefooter{}{}{}
	\runningfooter{}{}{}
\fi

\runningheadrule


\begin{document}
\pagestyle{headandfoot}   % exam class

%%%%%%%%%%%%%%%%%%%%%%%%%%%%%%%%%%%%%%%%%%%%%%%%%%%%%%
% title block
%%%%%%%%%%%%%%%%%%%%%%%%%%%%%%%%%%%%%%%%%%%%%%%%%%%%%%

\begin{center}
USAF ACADEMY \\
DEPARTMENT OF ELECTRICAL AND COMPUTER ENGINEERING \\
LESSON \labLesson\ LAB - First Flowgraphs \\
\textbf{Lab Packet} \\
\thisterm \\[4mm]
\end{center}

\vspace{-4pt}
\textbf{Name:} \rule[-2pt]{2.6in}{0.4pt} \hfill
\textbf{Section:} \rule[-2pt]{0.9in}{0.4pt} \hfill
\textbf{Date:} \rule[-2pt]{1.1in}{0.4pt}

\vspace{4pt}

\fullwidth{\textbf{Learning Objective 2.3: } I can install GNU Radio and use GNU Radio
Companion (GRC) to build, inspect, and run a basic flowgraph.
\\[2pt]\textbf{Learning Objective 2.4: } I can capture a real RF signal using an SDR
and interpret its time- and frequency-domain representations.}

\vspace{2pt}

This packet is the L\labLesson\ Lab, in three parts, and you turn it in once, complete,
before Lesson 6. Do \textbf{Part 1} on your own before the lab period: install
GNU Radio and work out the numbers your block settings fix. \textbf{Part 2} is
your own measurements and observations, with your own laptop and RTL-SDR,
recorded in class. Do \textbf{Part 3} on your own after the lab: explain what you
saw.

\textbf{Grading.} One submission, graded against Objectives 2.3 and 2.4 on the
course's 1 / 0 scale, with revision allowed under course policy. Paste screenshots
where a question asks for one.

\textbf{Receive only.} Nothing in this lab transmits. The RTL-SDR cannot transmit,
and the tone flowgraph never leaves your laptop.

The flowgraphs and the Python script are on the L5 lesson page:
\texttt{L05-tone.grc}, \texttt{L05-fm-spectrum.grc}, and \texttt{L05-psd.py}.

\begin{questions}

\fullwidth{\large\textbf{Part 1: Pre-Lab} \normalsize(individual, before the lab period)}

%%%%%%%%%%%%%%%%%%%%%%%%%%%%%%%%%%%%%%%%%%%%%%%%%%%
\question \textbf{Install and verify.} Install radioconda as the L5 lesson page
describes, plug in your RTL-SDR, and run the commands below from the radioconda
prompt.

\begin{parts}

	\begin{minipage}[t]{0.58\linewidth}
		\part Run \texttt{gnuradio-config-info --version} and record the version.
		\begin{solution}[0.5in]
		Any 3.10 release, for example 3.10.9.2. A 3.9 or earlier version means an
		old install is first on the path.
		\end{solution}
	\end{minipage}\hfill%
	\begin{minipage}[t]{0.37\linewidth}
		\raggedleft \vspace{12pt}
		version = \ansbox[1.3in]{3.10.x}
	\end{minipage}

	\begin{minipage}[t]{0.58\linewidth}
		\part Run \texttt{rtl\_test} and record the tuner it reports. Stop it with
		Ctrl-C after a few seconds.
		\begin{solution}[0.5in]
		Most course dongles report a Rafael Micro R820T or R828D tuner.
		\end{solution}
	\end{minipage}\hfill%
	\begin{minipage}[t]{0.37\linewidth}
		\raggedleft \vspace{12pt}
		tuner = \ansbox[1.3in]{R820T / R828D}
	\end{minipage}

	\part Open \texttt{L05-tone.grc} in GRC, run it, and paste a screenshot of the
	running flowgraph's window.
		\begin{solution}[1.6in]
		The window shows the frequency slider, the time sink with $I$ and $Q$ a quarter
		cycle apart, and the frequency sink with one spike at $+1$ kHz.
		\end{solution}

\end{parts}

%%%%%%%%%%%%%%%%%%%%%%%%%%%%%%%%%%%%%%%%%%%%%%%%%%%
\newpage
\question \textbf{Numbers from your settings.} Quick arithmetic from the block settings; you will check each on the screen at the bench.

\begin{parts}

	\part The tone flowgraph samples at $f_s = 32$ kHz. Work out where the frequency
	sink shows each tone.

	\begin{center}
	\renewcommand{\arraystretch}{1.25}
	\begin{tabular}{| l | K{0.68in} | K{0.68in} | K{0.68in} | K{0.68in} | K{0.68in} |}
	\hline
	Tone (kHz) & $1$ & $12$ & $20$ & $30$ & $-25$ \\\hline
	Expected (kHz) & \kv{$1$} & \kv{$12$} & \kv{$-12$} & \kv{$-2$} & \kv{$7$} \\\hline
	\end{tabular}
	\end{center}

		\begin{solution}[0.6in]
		Tones inside $\pm 16$ kHz appear where they are. Outside, subtract or add
		$32$ kHz until the result lands inside: $20 - 32 = -12$, $30 - 32 = -2$,
		$-25 + 32 = 7$.
		\end{solution}

	\begin{minipage}[t]{0.58\linewidth}
		\part The FM flowgraph samples at 2.4 MS/s centered on 100.1 MHz. Find the
		band the display shows.
		\begin{solution}[0.7in]
		\eq{ 100.1 \pm 1.2 = 98.9 \text{ to } 101.3\ \text{MHz} }
		\end{solution}
	\end{minipage}\hfill%
	\begin{minipage}[t]{0.37\linewidth}
		\raggedleft \vspace{12pt}
		band = \ansbox[1.5in]{98.9 to 101.3 MHz}
	\end{minipage}

	\begin{minipage}[t]{0.58\linewidth}
		\part Find the size of a 2 s recording at 2.4 MS/s.
		\begin{solution}[0.7in]
		\eq{ 2 \times 2.4\times10^{6} \times 8\ \text{bytes} = 38.4\ \text{MB} }
		\end{solution}
	\end{minipage}\hfill%
	\begin{minipage}[t]{0.37\linewidth}
		\raggedleft \vspace{12pt}
		\ansbox[1.3in]{38.4 MB}
	\end{minipage}

\end{parts}

%%%%%%%%%%%%%%%%%%%%%%%%%%%%%%%%%%%%%%%%%%%%%%%%%%%
\newpage
\fullwidth{\large\textbf{Part 2: Measurements} \normalsize(your own laptop and RTL-SDR, in class)}

\question \textbf{Build the tone flowgraph.} Start GRC with an empty canvas. Do not
open \texttt{L05-tone.grc}: build it yourself from the Options block, a
\texttt{samp\_rate} variable of 32000, a QT GUI Range named \texttt{freq} from
$-48000$ to $48000$, a complex Signal Source, a Throttle, a QT GUI Time Sink, and
a QT GUI Frequency Sink.

\begin{parts}

	\part Paste a screenshot of your flowgraph on the GRC canvas. Every connection
	must be black, not red.
		\begin{solution}[1.6in]
		Matches the L5 lesson figure. Every port is blue (complex).
		\end{solution}

	\part Set each tone with the slider and record where the spike appears.

	\begin{center}
	\renewcommand{\arraystretch}{1.25}
	\begin{tabular}{| l | K{0.68in} | K{0.68in} | K{0.68in} | K{0.68in} | K{0.68in} |}
	\hline
	Tone (kHz) & $1$ & $12$ & $20$ & $30$ & $-25$ \\\hline
	Measured (kHz) & \kv{$1$} & \kv{$12$} & \kv{$-12$} & \kv{$-2$} & \kv{$7$} \\\hline
	\end{tabular}
	\end{center}

		\begin{solution}[0.4in]
		The measurements match Question 2(a) exactly; aliasing of a computed tone
		is exact.
		\end{solution}

\end{parts}

%%%%%%%%%%%%%%%%%%%%%%%%%%%%%%%%%%%%%%%%%%%%%%%%%%%
\newpage
\question \textbf{The FM band.} Connect the antenna to the RTL-SDR and run
\texttt{L05-fm-spectrum.grc}.

\begin{parts}

	\part Tune across 88 to 108 MHz and record the three strongest stations: the
	frequency, the peak level on the display, and the approximate width.

	\begin{center}
	\renewcommand{\arraystretch}{1.25}
	\begin{tabular}{| l | K{1.3in} | K{1.3in} | K{1.3in} |}
	\hline
	& Frequency (MHz) & Peak (dB) & Width (kHz) \\\hline
	Station 1 & \kv{odd tenth} & \kv{varies} & \kv{$\approx 200$} \\\hline
	Station 2 & & & \\\hline
	Station 3 & & & \\\hline
	\end{tabular}
	\end{center}

		\begin{solution}[0.5in]
		Local stations vary. Every frequency should end in an odd tenth, and each
		station should be about 150 to 250 kHz wide.
		\end{solution}

	\part Find the DC spike. Tune so that a station sits on it, then tune 300 kHz
	away. Describe what moves and what does not.
		\begin{solution}[0.9in]
		The station moves across the display as you retune, and the spike stays at the
		center frequency. The spike belongs to the receiver, not to the air.
		\end{solution}

	\part With the strongest station on screen, set the gain to each value and
	record the noise floor away from any station. Mark the first setting where
	extra signals appear between stations.

	\begin{center}
	\renewcommand{\arraystretch}{1.25}
	\begin{tabular}{| l | K{0.7in} | K{0.7in} | K{0.7in} | K{0.7in} | K{0.7in} |}
	\hline
	Gain (dB) & $0$ & $20$ & $30$ & $40$ & $49$ \\\hline
	Floor (dB) & & & & & \\\hline
	Spurs? & & & & & \\\hline
	\end{tabular}
	\end{center}

		\begin{solution}[0.5in]
		The floor stays put at low gain, where the converter's own noise dominates, then
		rises once the gain lifts the antenna's noise above it. Spurs between stations
		usually appear in the 40 to 49 dB range near strong stations.
		\end{solution}

\end{parts}

%%%%%%%%%%%%%%%%%%%%%%%%%%%%%%%%%%%%%%%%%%%%%%%%%%%
\newpage
\question \textbf{Record and plot.} Center the strongest station slightly off the
DC spike and set the gain you chose in Question 4(c).

\begin{parts}

	\begin{minipage}[t]{0.58\linewidth}
		\part Enable the File Sink, run the flowgraph for about 2 s, stop it, and record
		the file size.
		\begin{solution}[0.6in]
		About 38 MB, matching Question 2(c); the exact size depends on how long it ran.
		\end{solution}
	\end{minipage}\hfill%
	\begin{minipage}[t]{0.37\linewidth}
		\raggedleft \vspace{12pt}
		size = \ansbox[1.3in]{$\approx 38$ MB}
	\end{minipage}

	\part Run the plotting script with your center frequency, and paste the plot:
	\\[2pt]\texttt{python L05-psd.py fm\_capture.cfile --fs 2.4e6 --fc <center Hz>}
		\begin{solution}[1.8in]
		The plot reproduces the frequency sink's view: the same stations at the same
		frequencies, the DC spike at the center, and the noise floor between them.
		\end{solution}

\end{parts}

%%%%%%%%%%%%%%%%%%%%%%%%%%%%%%%%%%%%%%%%%%%%%%%%%%%
\newpage
\fullwidth{\large\textbf{Part 3: Writeup} \normalsize(individual, after the lab)}

\question \textbf{Analysis.}

\begin{parts}

	\part Compare your measured aliases in Question 3(b) with your answers in
	Question 2(a). Explain why a 20 kHz tone and a $-12$ kHz tone produce the same
	samples at 32 kHz.
		\begin{solution}[1.2in]
		They match exactly. On the sampling wheel, a tone 32 kHz higher turns one
		extra whole turn between samples, which lands each sample on the same point.
		$20$ kHz and $-12$ kHz differ by exactly $32$ kHz, so their samples are
		identical.
		\end{solution}

	\part Using your table in Question 4(c), describe in words what happened as you
	raised the gain from 0 to 49 dB: to the noise floor, to how far the strongest
	station stood above it, and to the space between stations.
		\begin{solution}[1.5in]
		At low gain the floor stayed put, the converter's own noise, and the station
		rose with the gain, so it stood out more. Once the floor began to rise, the
		antenna's noise rose with the station, and the gap stopped growing. At the
		highest settings false signals appeared between stations: the strong station
		was overloading the converter.
		\end{solution}

	\part State the gain you would use to receive weak stations in this room, and
	justify it from your table in Question 4(c).
		\begin{solution}[1.0in]
		The setting where the floor has risen a few dB above its low-gain value but no
		spurs have appeared, typically 30 to 40 dB. Below it the converter's noise buries
		weak stations; above it strong stations overload the converter.
		\end{solution}

\end{parts}

\end{questions}

\vspace{1.0em}
\noindent\textbf{Documentation:} \rule[-2pt]{0.6\linewidth}{0.4pt}

\end{document}
